\section*{Введение}
\addcontentsline{toc}{section}{Введение}

Разведочный анализ данных (Exploratory Data Analysis, EDA) представляет
собой совокупность методов, предназначенных для первичного изучения
данных. Он позволяет получить представление о структуре исследуемого
набора данных, оценить основные статистические характеристики,
выявить выбросы и аномалии, а также обнаружить взаимосвязи между
переменными.
 
Проведение разведочного анализа является важным этапом
статистического исследования, поскольку результаты такого анализа
помогают выбрать подходящие методы дальнейшего моделирования и
сформулировать проверяемые статистические гипотезы.
 
\section*{\normalfont Описание выбранного датасета}

В работе используется датасет \textit{Wine Quality}, загруженный из файла \texttt{wine\_quality\_merged.csv}. Датасет объединяет два исходных набора данных — \textit{Red Wine Quality} и \textit{White Wine Quality} — и содержит информацию о физико-химических свойствах красных и белых вин, а также экспертные оценки их качества.
 
\begin{itemize}
    \item \textbf{Название датасета:} Wine Quality (Red and White Wine Quality).

    \item \textbf{Описание:} Датасет содержит измерения физико-химических характеристик образцов вина. Для каждого наблюдения указаны показатели кислотности, содержание сахара, хлоридов, диоксида серы, плотность, pH, содержание сульфатов и алкоголя, а также экспертная оценка качества вина и его тип (красное или белое).

    \item \textbf{Количество наблюдений:} 6497.

    \item \textbf{Количество признаков:} 13.

    \item \textbf{Тип признаков:} одиннадцать количественных непрерывных
    признака:
    \begin{itemize}
        \item \texttt{fixed acidity (фиксированная кислотность)};
        \item \texttt{volatile acidity (летучая кислотность)};
        \item \texttt{citric acid (лимонная кислота)};
        \item \texttt{residual sugar (остаточный сахар)}
        \item \texttt{chlorides (хлориды)};
        \item \texttt{free sulfur dioxide (свободный диоксид серы)};
        \item \texttt{total sulfur dioxide (общий диоксид серы)};
        \item \texttt{density (плотность)};
        \item \texttt{pH (кислотность)};
        \item \texttt{sulphates (сульфаты)};
        \item \texttt{alcohol (содержание алкоголя)}.
    \end{itemize}
    Один порядковый категориальный признак ---
    \texttt{quality} (оценка качества). Один номинальный категориальный признак --- \texttt{type} (тип вина: \texttt{red} или \texttt{white}).

    \item \textbf{Шкалы измерений:} количественные признаки измерены по шкале отношений, признак \texttt{quality} --- по порядковой шкале, признак \texttt{type} --- по номинальной шкале. Интервальная шкала в датасете не представлена.

    \item \textbf{Наличие пропусков:} пропущенные значения отсутствуют. Общее количество пропусков равно нулю. При этом в датасете обнаружено 1169 дублирующихся наблюдений.
\end{itemize}

\section*{\normalfont Методика исследования}

Исследование проводилось в несколько этапов с использованием языка
программирования Python и библиотек \texttt{pandas}, \texttt{numpy},
\texttt{matplotlib}, \texttt{seaborn} и \texttt{scipy}.
 
\begin{enumerate}
    \item На первом этапе был загружен датасет Wine Quality из файла \texttt{wine\_quality\_merged.csv}.
    Были определены размер датасета, количество наблюдений и признаков, а также типы данных отдельных переменных..

    \item На втором этапе была выполнена первичная проверка качества
    данных. Для каждого признака проверялось наличие пропущенных
    значений и дубликатов.

    \item На третьем этапе для количественных признаков были рассчитаны основные описательные статистики: среднее арифметическое, медиана, мода, стандартное отклонение, асимметрия и эксцесс. Сравнение среднего значения и медианы использовалось для предварительной оценки асимметрии распределений и анализа корректности вычисления среднего.

    \item На четвёртом этапе были построены графики для проверки
    фундаментальных предположений измерительного процесса. Использовались
    графики последовательности наблюдений, лаговые диаграммы,
    гистограммы и Q--Q-графики.

    \item График последовательности наблюдений использовался для выявления изменений среднего уровня и возможной упорядоченности данных. Лаговая диаграмма применялась для оценки зависимости между соседними наблюдениями.

    \item Гистограммы использовались для анализа формы распределений, выявления асимметрии и наличия нескольких групп значений. Q–Q-графики применялись для предварительной оценки соответствия распределений нормальному закону.

    \item На пятом этапе были исследованы распределения количественных признаков с помощью гистограмм с различным числом бинов, Box Plot, эмпирических функций распределения (CDF) и ядерной оценки плотности (KDE). Выполнена проверка потенциальных выбросов по правилу 1,5 · IQR.
    
    \item На шестом этапе был проведён анализ взаимосвязей между признаками. Использовались диаграммы рассеяния с линиями линейной регрессии, ковариационная матрица, матрица коэффициентов корреляции Пирсона и матрица диаграмм рассеяния. Дополнительно проверялось влияние потенциальных выбросов на корреляционные связи.
    
    \item На седьмом этапе был выполнен сравнительный анализ распределений количественных признаков по группам, заданным категориальной переменной type. Сравнение проводилось с помощью Box Plot по группам, эмпирических функций распределения и численного сопоставления медиан и разброса.
    
    \item На заключительном этапе результаты визуального и статистического анализа были сопоставлены между собой. На основании полученных результатов были сформулированы статистические выводы и гипотезы для дальнейшего моделирования.
\end{enumerate}
